\section{The minimal examples}\label{sec:examples}

Throughout, vertices are labelled $0,\dots,n-1$ as in the \texttt{graph6} strings given in the text and in the repository, so that every statement can be checked mechanically.

\subsection{Thirteen vertices: odd torsion}\label{sec:thirteen}

\begin{theorem}\label{thm:A}
Every $\Ztwo$-acyclic graph on at most $12$ vertices is strong I-contractible. On $13$ vertices there are exactly three $\Ztwo$-acyclic graphs that are not strong I-contractible, namely
\begin{center}\small
$\Gt=$ \gsix{L?qa``\{wttym\textasciitilde{}o}\ ($40$ edges),\qquad $\Gt+\{4,7\}$ ($41$),\qquad $\Gt+\{4,7\}+\{5,6\}$ ($42$).
\end{center}
$\Gt$ is the unique stuck graph on $13$ vertices and the other two form its $\mathrm{I}_4$-closure. For all three, $\widetilde H_*(\,\cdot\,;\ZZ)$ is $\ZZ/3$ in degree $1$ and zero otherwise; in particular none of them is I-contractible, and $\Gt$ is a flag triangulation of the Moore space $M(\ZZ/3,1)$ with $28$ triangles.
\end{theorem}

\begin{proof}
The first statement is the equality of the columns ``$\Ztwo$-acyclic'' and ``$\IS$'' in Table~\ref{tab:card} for $n\le 12$ (exact decision of $\IS$ for each $\Ztwo$-acyclic graph). For $n=13$, Theorem~\ref{thm:C} applies: the enumeration of all $2$-connected graphs with $4\le\delta$, $\Delta\le 11$ (Table~\ref{tab:runs}) found exactly one stuck graph, $\Gt$, and its $\mathrm{I}_4$-closure (Algorithm~\ref{alg:icert} without stopping rule) consists of the two graphs listed, both outside $\IS$ (every edge $uv$ with $\Nb(u,v)\in\IS$ was added; the three graphs are pairwise non-isomorphic since they have different numbers of edges). The homology is computed by Smith normal form (file \texttt{certificates/\allowbreak integer\_homology.txt}) and was cross-checked with \texttt{sympy} and Mathematica. Since graphs in $\I$ have contractible, hence $\ZZ$-acyclic, clique complexes, none of the three is in $\I$.
\end{proof}

\begin{figure}[t]
\centering
\IfFileExists{figures/fig01.pdf}{\includegraphics[width=\textwidth]{fig01.pdf}}{\fbox{\parbox{0.9\textwidth}{\centering\vspace{3cm}fig01 (pending)\vspace{3cm}}}}
\caption{The graph $\Gt$ on $13$ vertices. (a) Its clique complex cut open along the four edges lying in three triangles ($8\text{--}10$, $10\text{--}9$, $9\text{--}11$, $11\text{--}8$) is a disc with $28$ triangles and $21$ vertices whose boundary word is $(10\,9\,11\,8)^3$: the boundary winds three times around the $4$-cycle in the same direction. (b) The resulting space is the Moore space $M(\ZZ/3,1)$. (c) A plane drawing of $\Gt$ with vertex $12$ (degree $8$) at the centre, its neighbourhood $C_8$ around it and the $4$-cycle in red.}\label{fig:G13}
\end{figure}

The structure of $\Gt$ is shown in Fig.~\ref{fig:G13}. Its clique complex is pure $2$-dimensional; exactly four edges lie in three triangles and they form the induced $4$-cycle $C=8\text{--}10\text{--}9\text{--}11$. Cutting along $C$ (the only edges in more than two triangles) yields a disc whose boundary of length $12$ is mapped onto $C$ with the word $(10\,9\,11\,8)^3$, that is, with the pattern $a^3$: $\cl(\Gt)\simeq M(\ZZ/3,1)$. The automorphism group has order $8$, generated by $(0\,1\,3\,2)(4\,6\,7\,5)(8\,10\,9\,11)$ and $(0\,1)(2\,3)(5\,6)(8\,10)(9\,11)$; the degree sequence is $5^4\,6^4\,7^4\,8$. The open neighbourhoods of the nine vertices off $C$ are induced cycles ($C_6$ for $0,1,2,3$; $C_5$ for $4,5,6,7$; $C_8$ for $12$) and those of the four vertices of $C$ have seven vertices and eight edges (a cycle with a chord path), so no vertex is $\IS$-deletable; no common neighbourhood $\Nb(u,v)$ is $\Ztwo$-acyclic either (each is non-empty, since every edge lies in a triangle, and has at least two vertices and no cone point), so no edge is deletable---this is what ``stuck'' means.

\begin{remark}[Torsion in small flag complexes]
Adamaszek~\cite{Adamaszek2014} proved that $H_1(\cl(G);\ZZ)$ is torsion-free for every graph on at most $10$ vertices, that exactly four graphs on $11$ vertices have torsion in $H_1$, and that every graph on $12$ vertices with torsion in $H_1$ either is one of $363$ listed graphs or has a vertex $v$ such that $G-v$ is one of the four $11$-vertex graphs. All the listed complexes have torsion subgroup $\ZZ/2$ (they collapse to, or are homotopy equivalent to, $\mathbb{RP}^2$ or $\mathbb{RP}^2\vee S^1$); and if $G-v$ is one of the four $11$-vertex graphs, the Mayer--Vietoris sequence of $\cl(G)=\cl(G-v)\cup\mathrm{cone}(\cl(\Nb(v)))$ shows that the torsion of $H_1(G)$ is a quotient of $\ZZ/2$. Hence no flag complex on at most $12$ vertices has odd torsion in $H_1$, and Theorem~\ref{thm:A} shows that $13$ vertices suffice for $\ZZ/3$: the three graphs above are the smallest flag complexes with odd torsion in $H_1$. Our search covered only the $\Ztwo$-acyclic graphs on $13$ vertices, so we do not know whether other $13$-vertex flag complexes have odd torsion in $H_1$ (Question~\ref{q:torsion}).
\end{remark}

\begin{figure}[t]
\centering
\IfFileExists{figures/fig02.pdf}{\includegraphics[width=\textwidth]{fig02.pdf}}{\fbox{\parbox{0.9\textwidth}{\centering\vspace{3cm}fig02 (pending)\vspace{3cm}}}}
\caption{(a) The $15$-vertex flag dunce hat $\Dunce$ of~\cite[Fig.~2A]{DEFH2023} unfolded: a disc with $23$ vertices whose boundary of length $12$ runs along the $4$-cycle $1\text{--}2\text{--}3\text{--}4$ with the word $1\,2\,3\,4\;|\;1\,2\,3\,4\;|\;1\,4\,3\,2$, i.e.\ the pattern $a\,a\,a^{-1}$ (labels $1,\dots,15$ of~\cite{DEFH2023}; the \texttt{graph6} label $i$ corresponds to the label $(1,4,13,12,11,3,14,15,10,5,2,6,7,8,9)_i$, so this cycle is $0\text{--}10\text{--}5\text{--}1$ in \texttt{graph6} labels). (b) Reversing one of the three passes of the boundary of $\Gt$ (exchanging $10\leftrightarrow 11$ in one block) gives a dunce hat, but its $1$-skeleton contains the cliques $1\text{--}4\text{--}11$ and $2\text{--}5\text{--}10$ (purple) that are not faces of the disc, so the complex is not flag; the resulting graph is vertex-strong I-contractible. (c) The two identification patterns occurring among the stuck graphs of this paper.}\label{fig:dunce-vs-G13}
\end{figure}

\begin{proposition}\label{prop:no-flag-dunce-13}
There is no flag triangulation of the dunce hat on $13$ vertices whose $1$-skeleton is $\Ztwo$-acyclic and not strong I-contractible; in fact there is no graph on at most $13$ vertices in $\I\setminus\IS$.
\end{proposition}
\begin{proof}
By Theorem~\ref{thm:A} the only $\Ztwo$-acyclic graphs on $\le 13$ vertices outside $\IS$ are the three torsion graphs, none of which is contractible.
\end{proof}

It is instructive to see what goes wrong when one tries to turn $\Gt$ into a dunce hat (Fig.~\ref{fig:dunce-vs-G13}(b)): reversing the orientation of one block of the boundary word changes the pattern to $a\,a\,a^{-1}$ and the resulting complex is contractible, but two triples of pairwise adjacent vertices are no longer faces, so the complex is not the clique complex of its $1$-skeleton, and that $1$-skeleton (the graph \gsix{L?qa``\{wtlzM\textasciitilde{}o}, isomorphic to \gsix{LtDzQNKmpKXT`V}) contains the two copies $\{1,4,8,11\}$, $\{2,5,8,10\}$ of $K_4$ and is in $\IVS$.

\subsection{Fourteen vertices: the first examples of $\I\setminus\IS$}\label{sec:fourteen}

\begin{theorem}\label{thm:B}
There is no graph on at most $13$ vertices that is I-contractible but not strong I-contractible. Among $K_4$-free graphs on $14$ vertices there are exactly $51$ stuck graphs. Exactly three of them have trivial integer homology,
\begin{center}\small
$\Ga=$ \gsix{M?b@bFHNEk@\textasciitilde{}|k\}U?},\qquad
$\Gb=$ \gsix{M?B@thku?\textasciitilde{}Lylk\}[?},\\[2pt]
$\Gc=$ \gsix{M?r@d`ifbiFfo\}\textasciitilde{}o?},
\end{center}
with $44$, $45$ and $45$ edges and $31$, $32$ and $32$ triangles; each is I-contractible but not strong I-contractible. The remaining $48$ have $H_1(\,\cdot\,;\ZZ)\cong\ZZ/3$. Consequently the smallest graphs in $\I\setminus\IS$, and also the smallest $\ZZ$-acyclic graphs outside $\IS$, have exactly $14$ vertices; $\Ga$ has the minimum number of edges among the $K_4$-free ones.
\end{theorem}

\begin{proof}
The first statement is Proposition~\ref{prop:no-flag-dunce-13}; the same argument shows that no $\ZZ$-acyclic graph on at most $13$ vertices lies outside $\IS$ (such a graph would be $\Ztwo$-acyclic, hence one of the three torsion graphs, which are not $\ZZ$-acyclic), while $\Ga$ is $\ZZ$-acyclic and outside $\IS$. The enumeration of $K_4$-free $2$-connected graphs on $14$ vertices with $4\le\delta$, $\Delta\le 12$ (Table~\ref{tab:runs}; $\num{2000000}$ slices, all accounted for) returned $51$ stuck graphs, pairwise non-isomorphic by canonical labelling; their integer homology was computed by Smith normal form. For $\Ga$, the common neighbourhood of the non-adjacent vertices $0$ and $1$ is the induced path $8\text{--}12\text{--}5\text{--}10\text{--}13$, which is in $\IS$, so $\Ga+\{0,1\}$ is obtained from $\Ga$ by a transformation (I4); Algorithm~\ref{alg:IS} finds a sequence of $18$ strong deletions reducing $\Ga+\{0,1\}$ to $K_1$ (\ref{app:certificates}), so $\Ga+\{0,1\}\in\IS\subseteq\I$ and $\Ga\in\I$. Similarly $\Nb_{\Gb}(0,3)=7\text{--}13\text{--}9\text{--}12$ and $\Nb_{\Gc}(2,3)=6\text{--}13\text{--}7\text{--}10\text{--}11$ are induced paths and $\Gb+\{0,3\}$, $\Gc+\{2,3\}$ are in $\IS$ with $19$ and $18$ deletions. The three graphs are stuck, so not in $\IS$. Finally, a $K_4$-free graph in $\I\setminus\IS$ on $14$ vertices with fewer than $44$ edges would, by Lemma~\ref{lem:stuck}(b) (edge deletions preserve $K_4$-freeness), lie in the $\mathrm{I}_4$-closure of a stuck $K_4$-free graph with at most $43$ edges; all such stuck graphs have $H_1\cong\ZZ/3$, and (I4) additions preserve the homotopy type, a contradiction.
\end{proof}

\begin{figure}[p]
\centering
\IfFileExists{figures/fig03.pdf}{\includegraphics[width=\textwidth]{fig03.pdf}}{\fbox{\parbox{0.9\textwidth}{\centering\vspace{3cm}fig03 (pending)\vspace{3cm}}}}
\caption{$\Ga$ ($44$ edges, $31$ triangles, trivial automorphism group). (a) Unfolded disc with $24$ vertices; the boundary of length $15$ is identified with the $5$-cycle $5\text{--}10\text{--}13\text{--}9\text{--}12$ with the dunce-hat pattern $a\,a\,a^{-1}$ (boundary word $12\,5\,10\,13\,9\;|\;12\,5\,10\,13\,9\;|\;12\,9\,13\,10\,5$). (b) Plane drawing; the dashed edge $\{0,1\}$ is the (I4) addition certifying I-contractibility, and its common neighbourhood, the induced path $8\text{--}12\text{--}5\text{--}10\text{--}13$, is highlighted.}\label{fig:G14a}
\end{figure}
\begin{figure}[p]
\centering
\IfFileExists{figures/fig04.pdf}{\includegraphics[width=\textwidth]{fig04.pdf}}{\fbox{\parbox{0.9\textwidth}{\centering\vspace{3cm}fig04 (pending)\vspace{3cm}}}}
\caption{$\Gb$ ($45$ edges, $32$ triangles, $|\mathrm{Aut}|=2$). The boundary of the unfolded disc ($25$ vertices, length $16$) is identified with the $4$-cycle $5\text{--}7\text{--}13\text{--}8$ with pattern $a\,a\,a^{-1}$, and additionally runs back and forth along the pendant edge $13\text{--}9$, which lies in four triangles. The (I4) edge is $\{0,3\}$ with $\Nb(0,3)=7\text{--}13\text{--}9\text{--}12$.}\label{fig:G14b}
\end{figure}
\begin{figure}[p]
\centering
\IfFileExists{figures/fig05.pdf}{\includegraphics[width=\textwidth]{fig05.pdf}}{\fbox{\parbox{0.9\textwidth}{\centering\vspace{3cm}fig05 (pending)\vspace{3cm}}}}
\caption{$\Gc$ ($45$ edges, $32$ triangles, $|\mathrm{Aut}|=2$). The boundary of the unfolded disc ($26$ vertices, length $18$) is identified with the $6$-cycle $0\text{--}13\text{--}7\text{--}10\text{--}11\text{--}9$ with pattern $a\,a\,a^{-1}$. The (I4) edge is $\{2,3\}$ with $\Nb(2,3)=6\text{--}13\text{--}7\text{--}10\text{--}11$.}\label{fig:G14c}
\end{figure}

Figures~\ref{fig:G14a}--\ref{fig:G14c} show the three graphs. All three clique complexes are pure $2$-dimensional flag triangulations of a disc with its boundary identified along a short induced cycle (with, for $\Gb$, a pendant edge traversed back and forth) with the dunce-hat pattern $a\,a\,a^{-1}$ (Table~\ref{tab:compare}); they are therefore homotopy equivalent to the dunce hat and contractible but, like the dunce hat, presumably not collapsible (see Section~\ref{sec:open}). Compared with $\Dunce$ they use a longer identification cycle ($5$, $4$ and $6$ instead of $4$), fewer edges ($44$--$45$ instead of $46$) and a more homogeneous degree sequence (degrees between $5$ and $8$ instead of between $4$ and $10$). In $\Dunce$, $\Ga$ and $\Gc$ the vertices outside the identification cycle have an induced cycle as open neighbourhood (so their links are circles: these are interior points of the surface), and the vertices on the cycle have a neighbourhood with one more independent cycle; in $\Gb$ this holds for all vertices except the end $9$ of the pendant edge $13\text{--}9$, which lies in four triangles.

\begin{table}[t]
\centering\small
\caption{The three minimal examples and the dunce hat of~\cite{DEFH2023}. $|E_{\ge 3}|$ = number of edges in three or more triangles; disc = number of vertices of the unfolded disc; $k$ = length of the identification cycle.}\label{tab:compare}
\begin{tabular}{@{}lrrrlrrrrl@{}}
\toprule
 & $n$ & $m$ & triangles & degrees & $|E_{\ge 3}|$ & $|\mathrm{Aut}|$ & disc & $k$ & pattern \\
\midrule
$\Dunce$ & 15 & 46 & 32 & $4^3\,5^2\,6^7\,9^2\,10$ & 4 & 2 & 23 & 4 & $a\,a\,a^{-1}$ \\
$\Ga$ & 14 & 44 & 31 & $5^3\,6^6\,7^3\,8^2$ & 5 & 1 & 24 & 5 & $a\,a\,a^{-1}$ \\
$\Gb$ & 14 & 45 & 32 & $5^2\,6^5\,7^6\,8$ & 5 (one in four triangles) & 2 & 25 & 4 & $a\,a\,a^{-1}$ \\
$\Gc$ & 14 & 45 & 32 & $6^9\,7^4\,8$ & 6 & 2 & 26 & 6 & $a\,a\,a^{-1}$ \\
\bottomrule
\end{tabular}
\end{table}

\begin{figure}[p]
\centering
\IfFileExists{figures/fig06.pdf}{\includegraphics[height=0.9\textheight]{fig06.pdf}}{\fbox{\parbox{0.9\textwidth}{\centering\vspace{3cm}fig06 (pending)\vspace{3cm}}}}
\caption{Certificate that $\Ga+\{0,1\}\in\IS$: the $18$ strong deletions found by Algorithm~\ref{alg:IS} (deleted element in red, its neighbourhood or common neighbourhood in amber). The sequence is listed in \ref{app:certificates} and reproduced by \texttt{tools/strong\_contractible.py --trace}.}\label{fig:cert}
\end{figure}

The $48$ torsion examples are listed in \ref{app:stuck51} (Table~\ref{tab:stuck51}) and shown in Fig.~\ref{fig:mosaic}. They have $43$ ($25$ graphs), $44$ ($18$), $45$ ($4$) or $47$ ($1$) edges. For $45$ of them the unfolding is a disc whose boundary is identified with a cycle of length $4$ ($25$ graphs), $5$ ($18$) or $6$ ($2$) with the pattern $a^3$, as for $\Gt$; two unfold to an annulus and one (the graph with $47$ edges, whose edges in three triangles form two disjoint $4$-cycles) to a disconnected surface, so their $\ZZ/3$ does not come from a single disc identification.

\begin{figure}[p]
\centering
\IfFileExists{figures/fig07.pdf}{\includegraphics[height=0.92\textheight]{fig07.pdf}}{\fbox{\parbox{0.9\textwidth}{\centering\vspace{3cm}fig07 (pending)\vspace{3cm}}}}
\caption{The $51$ stuck $\Ztwo$-acyclic $K_4$-free graphs on $14$ vertices: the three contractible ones ($\I\setminus\IS$) first, then the $48$ with $H_1\cong\ZZ/3$, ordered by number of edges.}\label{fig:mosaic}
\end{figure}

\paragraph{The $\mathrm{I}_4$-closures} By Theorem~\ref{thm:C}(ii), every $K_4$-free graph in $\I\setminus\IS$ on $14$ vertices lies in the $\mathrm{I}_4$-closure of one of the $51$ stuck graphs (the closure itself may leave the $K_4$-free class). Computing all $51$ closures gives $\num{2271}$ pairwise non-isomorphic $\Ztwo$-acyclic graphs on $14$ vertices outside $\IS$, of which $29$ are contractible, hence in $\I\setminus\IS$, and $\num{2242}$ have $H_1\cong\ZZ/3$ (file \texttt{data/acyclic\_not\_\allowbreak strong\_14\_closure.tsv} in the repository); $8$ of the $29$ are $K_4$-free, with $44$--$47$ edges. Hence the $K_4$-free graphs on $14$ vertices in $\I\setminus\IS$ are exactly these $8$ graphs. (The closures do not exhaust the $\Ztwo$-acyclic graphs outside $\IS$ on $14$ vertices: adding to one of the three $13$-vertex torsion graphs a new vertex whose neighbourhood is in $\IS$ produces further graphs with $H_1\cong\ZZ/3$ outside $\IS$, which are not stuck and not in these closures. Lemma~\ref{lem:stuck}(a) does not apply at $n=14$ since $n_0=13$.)

\subsection{Relation with the dunce hat of~\cite{DEFH2023}}

The graph $\Dunce$ (Fig.~\ref{fig:dunce-vs-G13}(a)) is also stuck, with $46$ edges, $32$ triangles and the identification pattern $a\,a\,a^{-1}$ along the $4$-cycle (the $4$-cycle $1\text{--}2\text{--}3\text{--}4$ in the labels of~\cite[Fig.~2A]{DEFH2023}; the label map is recorded in \texttt{tools/dunce\_hat.py}). It belongs to the same family as $\Ga,\Gb,\Gc$: flag discs with identified boundary. It is not minimal because the same identification can be realised with one vertex fewer by letting the boundary wind around a $5$- or $6$-cycle, which spreads the three boundary passes over more edges and allows a more uniform (and smaller) vertex set. Adding the edge $\{9,12\}$ to $\Dunce$ (edge $\{5,7\}$ of~\cite[Fig.~2B]{DEFH2023}) produces a graph in $\IS$ with a $16$-step deletion sequence (\ref{app:certificates}), exactly as adding $\{0,1\}$ to $\Ga$ does.
